%% Macros de dados do documento
\autor{PEDRO HENRIQUE RIGUETE DE OLIVEIRA MARTINS}
\titulo{ Gestão Operacional e Controle de Lanchonetes: Desenvolvimento de um Sistema Informatizado para Automação de Processos}
\instituicao{Instituto federal do espírito santo}
\curso{Curso de Análise e Desenvolvimento de Sistemas}
\local{ALEGRE}
\data{2026}


\tipotrabalho{Trabalho de Graduação}
\preambulo{Trabalho de Conclusão de Curso apresentado à Coordenadoria do Curso de Análise e Desenvolvimento de Sistemas do
Instituto Federal de Educação, Ciência e Tecnologia do Espírito Santo, como requisito parcial para a obtenção do título de Analista em Desenvolvimento de Sistemas.}

\orientador[Orientador][Prof. Dr.]{Carlos Alexandre }[Siqueira da Silva]{Instituto Federal do Espírito Santo}


\examinadori{Profa. Dra. Fulana de Tal}{Instituto Federal do Espírito Santo}{Examinadora}
\examinadorii{Prof. Dr. Cicrano de Tal}{Instituto Federal do Espírito Santo}{Examinador}

\approvaldate{01}{Agosto}{2026}

% Palavras-chaves
\palavraschave{Palavra Chave 1}[Palavra Chave 2][Palavra Chave 3][Palavra Chave 4][Palavra Chave 5]
\keywords{keyword 1}[keyword 2][keyword 3][keyword 4][keyword 5]

% Ficha catalográfica
\fichabiblioteca{
    Dados Internacionais de Catalogação-na-Publicação (CIP) \\
    (Biblioteca Nilo Peçanha do Instituto Federal do Espírito Santo)
}
\fichaautor{Elaborada por XXXXXXXXXXXXXXXXXXXXX – CRB-X/ES - XXX}
\fichatags{
    1. Redes neurais (Computação).  2. Expressão facial – Processamento de dados. 3. Sistemas de reconhecimento de padrões. 4. Processamento de imagens – Técnicas digitais. 5. Percepção de padrões. 6. Engenharia Elétrica. I. XXXXXXXX, XXXXXXXXXX XXXXXXXXXX.  II. Instituto Federal do Espírito Santo. III. Título.
}

\cutter{X999y}
\cdd{000.00}
\selectlanguage{brazil}

% Arial (para usar times-new-roman, comente as três linhas abaixo)
\usepackage{helvet}
\renewcommand{\familydefault}{\sfdefault}
\renewcommand{\ABNTEXchapterfont}{\bfseries}
% Comandos para revisar o texto
\usepackage{easyReview} 


%% Elementos opcionais
\editardedicatoria{\input{pre_textuais/dedicatoria}}
\editaragradecimentos{\input{pre_textuais/agradecimentos}}
\editarepigrafe{\input{pre_textuais/epigrafe}}


%% Listas [Siglas e Simbolos]
\editarlistasiglas{\begin{siglas}
    \item[Ifes] Instituto Federal do Espírito Santo
    \item[PHP] Personal Home Page
    \item[mvc] model-view-controller 
    \item[RBAC]Role-Based Access Control ou Controle de Acesso Baseado em Funções
    \item[RNF] Requisitos Não Funcionais
    \item[RF] Requisitos Funcionais
    \item[SQL] Structured Query Language (Linguagem de Consulta Estruturada)
    \item[SGBD] Sistema de Gerenciamento de Banco de Dados
\end{siglas}}
\editarlistasimbolos{\input{pre_textuais/simbolos}}